\section{The coarse finite-cocycle lift on the cyclotomic field}

% BEGIN B81L1
\paragraph{B81L1. Intersection with the actual kernel lattice.}
Work at $F=K_0(\mu_p)$ on the case B layer of A81R2, with $p$ odd,
$A=\prod_i O[[X]]$, $Q=\Phi_p(1+X)$, and $u=\chi(\gamma)$,
$v_p(u-1)=1$. Let $N_j=N(T_j)$ for an exact sequence of crystalline
lattices $T_0\to T_1\xrightarrow{f} T_2$. Assume the middle jumps are in
$[-h,h]$; so are both outer sets by strictness (H80M1, J80H1).
We claim $D(T_0)\cap N_1=N_0$ inside $D(T_1)$.

Recall why one can take kernels in these coordinates. In bases of
the three \'etale modules, FET22--FET25 identify the maps after
extension to $W_0=O\otimes W(C^\flat)$ with those of the original
lattices tensored with $W_0$. This base change is faithfully flat
on the \'etale coefficient ring: modulo the uniformizer of $O$ it
is, componentwise, extension of the Laurent series field into
copies of the tilt field. The base components are complete DVRs
with this uniformizer (successive lifting of a residue-field inverse
inverts every element of valuation zero); the target is torsion-free,
and its reduction is nonzero on every component. Thus exactness
and injection descend. For an element of $D(T_0)$ mapping into $N_1$, its image lies in
$N_1\cap\ker f$. A81R10 identifies this intersection with the image
of $N_0$ on each letter. By injectivity of the \'etale map its
preimage is therefore in $N_0$. Conversely functoriality carries
$N_0$ into $N_1$. In particular both $Q^{-h}N_1$ and $XN_1$
intersect the kernel in the corresponding modules for $N_0$:
multiply by $Q^h$ for the first assertion, and for the second
write the vector as $Xn$ and use injectivity of $X$ on the target.
These statements apply also to the crystalline extensions used below.
% END B81L1

% BEGIN B81L2
\paragraph{B81L2. A bounded Herr pair gives a crystalline extension.}
Call $(y,z)$ a bounded pair for $T$ if $Q^hy\in N(T)$,
$z\in XN(T)$ and $(\gamma-1)y=(\varphi-1)z$. It gives an integral
finite class. We construct the comparison for its extension to
avoid assuming that this extension is crystalline.
Write $j$ for the Witt embedding and choose $b$ in $W_0\otimes T$
with $(\varphi-1)b=j(y)$ as in LWR5. Let $c$ be the actual
cocycle constructed there: $c(g)=d(g)-(g-1)b$ inside the Witt
tensor, where $d$ is zero on the cyclotomic kernel and
$d(\gamma)=j(z)$ on a generator lift. Give $T'=T\oplus O e$
the action $g e=e+c(g)$. In its Witt tensor the vector $v=e+b$
is fixed by the kernel, $(\varphi-1)v=j(y)$, and the generator
satisfies $(\gamma-1)v=j(z)$. The columns of a Wach basis of
$N(T)$ together with $v$ form a Witt basis: its change of basis
is triangular with the last diagonal entry one. Let $N'$ be their
$A$-span. It is free. The generator and its inverse preserve $N'$
and act identically modulo $X$ since $z\in XN(T)$. These facts
hold on the full cyclotomic group by density: $N'$ and $XN'$ are
compact images of finite powers of $A$ for its $(\lambda,X)$-adic
topology in the Hausdorff Witt-coordinate topology, with continuous
substitution as in LWR3--LWR5. The kernel fixes each basis column.

The Frobenius matrix in these columns is triangular by blocks with
last diagonal entry one, invertible after inverting $Q$, and both
it and its inverse have bounded powers of $Q$ as denominators.
This follows for its old block from W79H4, for its new upper column
from $Q^hy\in N(T)$, and then for the inverse by block inversion.
These verified matrix properties suffice for the following construction.
L78W2--L78W3 clear powers of $X$ in this Witt basis and its inverse;
apply the map from $A_0$ to $L\otimes B_{\rm cris}$ and divide by
these powers there using L78W4. The two resulting matrices remain
inverses. In particular these columns can be evaluated in the
crystalline tensor space of $T'[1/p]$ before its invariant dimension
has been determined. For the gamma matrix $G$, with $G(0)=1$,
L78W9 solves $G\gamma(U)=U$ and $U(0)=1$ by dividing the $j$th
coefficient equation by $1-u^j$, $j>0$. The commuting Frobenius
matrix and its inverse have just been shown to have finite $Q$-poles,
so the continuation and evaluation in L78W10 apply to $U$ and its
inverse. At every group element uniqueness of the normalized solution
intertwines this matrix with its Galois transform. Multiplying the
columns by the evaluated $U$ thus gives a $B_{\rm cris}$-basis of
invariants, with coefficients on the same finite extension as in
W79H1. Restricting the matrices along each coefficient idempotent,
or taking all coefficient coordinates, gives the required crystalline
comparison. In particular the trivial quotient has a fixed crystalline
lift: project the invariant basis, and use that the invariant scalars
are $L\otimes K_0$ to lift the quotient generator componentwise.
Subtract this lift from $e$ to obtain a crystalline primitive for $c$.
This proves the claimed finiteness from the bounded pair itself.
% END B81L2

% BEGIN B81L3
\paragraph{B81L3. The chosen finite class has a rationally liftable bounded pair.}
Write $b_W=w(4h)$ and use the original Wach cokernel bound of
A81R8--A81R9 on the cyclotomic field. For a target integral
finite cocycle $c$ form $\widetilde T_2=T_2\oplus Oe$ using $c$.
The rational extension is crystalline: the assumed crystalline
primitive yields an invariant lift of $e$ over $B_{\rm cris}$,
which together with a comparison basis in $V_2$ supplies a basis.
By strictness the extra graded piece has degree zero. In
$N(\widetilde T_2)$ lift $p^{b_W}$ times the fixed generator of
$N(O)$ to $v_2$. The differences
$y=(\varphi-1)v_2$, $z=(\gamma-1)v_2$ lie in the kernel \'etale
module. They give a bounded pair by the lower-height bound, trivial
action modulo $X$, and B81L1 on the kernel (also divide by $X$ in
the second coordinate). In the Witt tensor write $v_2=p^{b_W}e+b$.
Its two difference equations and invariance on the cyclotomic kernel
are the defining rules in LWR5, as checked also in A81R11. Accordingly its Herr class is
$p^{b_W}[c]$, with this sign.

The rational map $H^1_f(F,V_1)\to H^1_f(F,V_2)$ is surjective.
Indeed the three crystalline spaces are exact (J80H1), so
$H^1(F,B_{\rm cris}\otimes V_0)\to H^1(F,B_{\rm cris}\otimes V_1)$
is injective by the long exact sequence. The ordinary exact
sequence through $H^1(F,V_2)$ consequently restricts to an exact
sequence through $H_f^1(F,V_1)$, retaining the three $H^0$ terms.
Applying the crystalline dimension formula and strictness of the
filtered sequence shows that the image in $H_f^1(F,V_2)$ already
has its whole dimension: the dimensions of the three quotients
by $\mathrm{Fil}^0$ add in the short exact sequence and the $H^0$
terms cancel. Hence there is a rational crystalline extension
$\widetilde V_1\to\widetilde V_2$ inducing $f:T_1\to T_2$ on
kernels and the identity on quotient lines, after adjusting a
rational lift by a coboundary.

Choose a lattice $\widetilde T_1=T_1\oplus O p^l e_1$ in it with
$l$ large: compactness clears the values of the rational extension
cocycle, and the image of the added generator is integral in
$\widetilde T_2$. Its image is a lattice $I$ of finite index in
$\widetilde T_2$, containing $p^l\widetilde T_2$ (the kernel map
is onto). The integral maps $\widetilde T_2\xrightarrow{p^l}I$
and $I\hookrightarrow\widetilde T_2$ compose to multiplication
by $p^l$. Apply the Wach cokernel bound to
$\widetilde T_1\twoheadrightarrow I$ and the Wach functor to these
maps. For some $k$ this gives $v_1\in N(\widetilde T_1)$ lying
above $p^k v_2$. Its image in the \'etale quotient line is constant,
since the map on quotient lines is rationally the identity.
B81L1 and the same two difference calculations give a bounded
pair in $D(T_1)$ lying exactly above $(p^k y,p^k z)$. The size
of $k$ will have no effect on the final bound.
% END B81L3

% BEGIN B81L4
\paragraph{B81L4. The polynomial and the two integral regulators on the kernel.}
Use $N=N_0$, $D=D(T_0)$, and write $U_1$ for the set of all
jumps of the source $V_1$ over the embeddings of the cyclotomic
field. Use the set $J$ of A81R2, so that it contains each nonzero
$a+j$ with $a\in U_1$ and $1\le j\le h$. Put
\[
 P(B)=\prod_{j\in J}(B-u^{-j}),\qquad
 C=\prod_{i=1}^{h}(u^i-1),\qquad N_Q=Q^{-h}N.
\]
Then $(\gamma-1)P(\gamma)N_Q\subset N$. Indeed the quotient
$N_Q/N$ is free over $A/(Q^h)$ in a Wach basis, up to the
trivialization by $Q^h$. It is $p$-torsion free; Weierstrass
division for the distinguished polynomial $Q^h$ makes each
coefficient quotient finite free over $O$. After inverting $p$
this quotient is the product of its Taylor expansions of order
$h$ at the distinct primitive roots (Chinese remaindering).
In the coordinate $1+X=\zeta\exp(t')$, $Q$ is $t'$ times a
unit, and $\gamma$ acts on $t'$ by $u t'$. W79H2--W79H3 give
a filtered basis $b_a$ of invariant periods, indexed with
multiplicities by the jumps at this branch, in which the local
lattice has basis $(t')^{-a}b_a$. To verify the last description,
the degree-$k$ constant coefficients in the graded expansion
lattice consist precisely of $\mathrm{Fil}^{-k}$ by W79H2;
interpolation in W79H3 separates the degrees. A basis adapted
to this filtration consequently gives the displayed powers.
Thus in $Q^{-h}N/N$ the eigenvalues at this branch are
$u^{-a-j}$ for $1\le j\le h$, with no Jordan contribution.
These $a$ are jumps of $V_1$ by strictness. The claimed
polynomial kills every branch, including eigenvalue one, and
hence kills the integral quotient.

Two more calculations take place over $A$ without inverting $p$.
One has
\[
 C XN\subset(\gamma-1)XN+X^{h+1}N,
 \qquad (1-\varphi)X^{h+1}N=X^{h+1}N.
\]
For the first, $\gamma$ acts on $X^iN/X^{i+1}N$ by $u^i$.
Thus $\prod_{i=1}^h(\gamma-u^i)$ carries $XN$ into
$X^{h+1}N$. Divide this polynomial minus its value at one by
$B-1$ and absorb its sign to give the assertion. For the second,
$Q^h\varphi N\subset N$ gives
$\varphi(X^{h+1}N)\subset QX^{h+1}N$. The ideal $(p,X)$ is
Frobenius-stable and contains $Q$. Iteration therefore pushes
into $(p,X)^kX^{h+1}N$. The geometric series for $(1-\varphi)^{-1}$
converges in the complete power-series module. Finally both $p$
and $X$ are nonzerodivisors on $N_Q/N$. For $X$ this follows
from the same finite free coefficient quotient: after passing to
$L$, its roots $\zeta-1$ are nonzero. In particular
$N\cap XN_Q=XN$.
% END B81L4

% BEGIN B81L5
\paragraph{B81L5. Correcting the integral defect with the prescribed exponent.}
In the kernel module set $d(Y,Z)=(\gamma-1)Y-(\varphi-1)Z$.
The admissible inputs for $d$ form $N_Q\oplus XN$.
Both summands are $\gamma$-stable. Indeed $\gamma(Q)/Q$ is an
integral unit of constant coefficient one, since the substitution
permutes the primitive roots and $Q(0)=p$. It follows that
$(\gamma-1)N_Q\subset XN_Q$, and
$(\varphi-1)XN\subset XN_Q$ by the height bound. Thus
the defect group can be computed inside $XN_Q$.

Suppose $\omega\in XN_Q$ and, for some $k$, one has
$p^k\omega=d(Y,Z)$ with $Y\in N_Q$, $Z\in XN$. Write
$CZ=(\gamma-1)x+z_1$ with $x\in XN$ and
$z_1\in X^{h+1}N$ using B81L4. Commutativity gives
\[
 Cp^k\omega=(\gamma-1)(CY-(\varphi-1)x)-(\varphi-1)z_1.
\]
The first parenthesis is in $N_Q$; the last term is in $N$.
Apply $P(\gamma)$: its product with $\gamma-1$ kills $N_Q/N$,
so $Cp^kP(\gamma)\omega\in N$. Since $C\ne0$ and $N_Q/N$ is
$p$-torsion free, $P(\gamma)\omega\in N$. It is also in
$XN_Q$, hence lies in $XN$. The two regulators in B81L4 now give
$CP(\gamma)\omega\in d(N_Q,XN)$: the term in $X^{h+1}N$
is a $(1-\varphi)$-image there. The polynomial
$P(B)-P(1)$ is a multiple of $B-1$ over $\mathbb Z_p$, and
$\omega\in N_Q$. Subtracting its contribution gives
$CP(1)\omega\in d(N_Q,XN)$. The valuation of $CP(1)$ is exactly
\[
                 d_U=w(h)+\sum_{j\in J}(1+v_p(j)).
\]
Indeed $0\notin J$ and $v_p(u^j-1)=1+v_p(j)$ for every nonzero
integer $j$. A $p$-adic unit rescales $CP(1)$ to $p^{d_U}$.

For the pair $(y,z)$ of B81L3, lift $p^{b_W}Q^h y$ and
$p^{b_W} z/X$ integrally to $N_1$ by the Wach cokernel estimate.
Divide the first lift by $Q^h$ and multiply the second by $X$ to
obtain $(Y,Z)$ above $p^{b_W}(y,z)$. The element $d(Y,Z)$ lies
in the \'etale kernel, and is in $XQ^{-h}N_1$ by the calculation
above. B81L1 and injectivity of $X$ put it in $XN_Q$. This
defect satisfies the torsion hypothesis just treated: subtract
from $p^k(Y,Z)$ the $p^{b_W}$-multiple of B81L3's bounded
cocycle over $p^k(y,z)$. The differences are in the kernel
$N_Q\oplus XN$ by B81L1, and their differential is $p^kd(Y,Z)$.
Hence one can subtract from $p^{d_U}(Y,Z)$ a kernel element
of $N_Q\oplus XN$ whose differential is $p^{d_U}d(Y,Z)$. The result is a cocycle
pair on $T_1$ mapping to $p^{d_U+b_W}(y,z)$ in both coordinates.

By B81L2 the lifted pair has integral finite class. Naturality of
LWR5 follows with its explicit primitive by applying the lattice
map to all the Witt-coordinate equations. Its class maps to
$p^{d_U+2b_W}[c]$ by B81L3. This is exactly the case B exponent
$E$ of A81R2 at this member and on this restricted map. The
prime-to-$p$ corestriction in A81R2 gives the original coarse
lifting test at the two ramified primes, on every tested row.
% END B81L5
